% ===== roadmap of this week =====
\begin{frame}{이번 주차 개요}
  \textbf{이 주차를 마치면}
  \begin{itemize}
    \item 결정 트리의 구조를 설명하고, 지니불순도 감소량(정보이득)을 계산해
      ``좋은 질문''을 고른다
    \item 랜덤 포레스트의 두 가지 무작위성이 분산을 줄이는 이유를 설명하고,
      OOB 점수로 검증 세트 없이 일반화 성능을 예측한다
    \item 랜덤 포레스트(독립적·다수결)와 GBDT(순차적·잔차 추격)를 구분하고,
      GBDT의 트리 개수는 early stopping으로 정해야 하는 이유를 설명한다
    \item SHAP으로 예측 하나를 특징별 기여로 분해하고, 모든 SHAP값의 합이
      정확히 예측값을 재현하는 가산성을 확인한다
  \end{itemize}
  \vskip0.6em
  \begin{itemize}
    \item \kb{7.1} — 트리를 ``질문''으로 세우고, 지니불순도·정보이득·U자 곡선을 다룬다
    \item \kb{7.2} — 트리를 \textbf{병렬}로 합치는 랜덤 포레스트(독립 + 다수결)
    \item \kb{7.3} — 트리를 \textbf{순서}로 합치는 GBDT(잔차를 고쳐가는 보정)와 그 설명(SHAP)
  \end{itemize}
  \vskip0.6em
  \begin{block}{세 블록을 관통하는 줄기}
    \textbf{합치는 방식이 만든다.} 독립적으로 합치면(랜덤 포레스트) 흔들림만 줄고,
    순차적으로 합치면(GBDT) 정확도는 더 오르지만 과적합의 위험도 함께 온다.
    정확도의 대가로 잃은 ``읽을 수 있음''은 SHAP(게임이론)으로 되찾는다.
  \end{block}
\end{frame}
